\documentclass[11pt]{article}
\usepackage{icmlko}
\begin{document}

\title{자가 여섯인데 물음은 넷 반이다 --- 그리고 붕괴한 방향은 하나뿐이다}
\author{월드모델 실험실}
\date{노트 614 · 2026년 8월 5일}
\maketitle

\begin{abstract}
노트 $613$ 이 여섯 자의 표본 SD 를 다 냈지만 \textbf{자끼리의 상관은 안
쟀다}. ``다섯 자가 다 통과''가 얼마나 센 말인지는 \emph{독립인 자가 몇
개냐}에 달렸다. 한 부트스트랩 되풀이 안에서 여섯을 같이 계산해 결합
분포를 냈다($3{,}200$ 점).

\begin{center}
\begin{tabular}{lr}
\hline
 & 표본 성분 상관 \\
\hline
판 $\leftrightarrow$ 날짜 통제 판 & $\mathbf{+0.88}$ \\
판 $\leftrightarrow$ 무리 안 판 & $+0.11$ \\
판 계열 $\leftrightarrow$ 집 밖 셋 & 최대 $|0.05|$ \\
\hline
\textbf{유효 독립 자 수} & $\mathbf{4.73 / 6}$ \\
\hline
\end{tabular}
\end{center}

\textbf{중복은 한 쌍뿐이다.} 여섯 자의 고윳값이 $1.91 \cdot 1.08 \cdot 1.02
\cdot 0.99 \cdot 0.88 \cdot \mathbf{0.12}$ 로, \textbf{붕괴한 방향이 정확히
하나}(판과 날짜 통제 판의 차)다. 관문은 생각보다 세다.
\end{abstract}

\begin{figure}[t]
\centering
\includegraphics[width=\textwidth]{figs/f1.pdf}
\caption{\textbf{가} --- 표본 성분(믿을 수 있다).
\textbf{나} --- 씨앗 성분(못 믿는다).
\textbf{다} --- 붕괴한 방향은 하나뿐이다.}
\end{figure}

\section{같은 행을 쓰는데 독립인 자}

무리 안 판은 판과 \textbf{똑같은 $3{,}369$ 행}을 쓰는데 표본 상관이
$+0.11$ 이다. 그럴 수 있는 이유는 \textbf{묻는 것이 다르기} 때문이다 ---
판은 도메인 전체의 순위를 묻고, 무리 안 판은 \emph{무리(\texttt{grp})
안에서의} 순위를 묻는다. 행이 같아도 그 두 양이 흔들리는 방향은 거의
겹치지 않는다.

\begin{quote}
\textbf{자의 독립성은 자료가 아니라 물음이 정한다.} 이 판은 두 해 동안
``집 밖 짝이 독립이다''를 \emph{다른 자료}라서 그렇다고 읽어 왔다. 실제로는
\textbf{같은 자료에서도 다른 물음이면 독립}이고, 정작 겹친 것은 같은
자료에 같은 물음을 조금 다르게 던진 판과 날짜 통제 판이었다.
\end{quote}

\section{사전 등록 ① 은 아슬하게 빗나갔다}

``판 $\leftrightarrow$ 날짜 통제 판 표본 상관이 $0.9$ 위''로 미리 적었고
\textbf{$+0.88$} 이 나왔다. \textbf{빗나갔다} --- 실질은 바뀌지 않지만
그대로 적는다(노트 $578$: 사전 등록은 셋을 다 지킨다).

사전 등록 ②는 섰다 --- 판 계열과 집 밖 셋의 표본 상관이 최대 $|0.05|$ 다.
\textbf{두 자료가 겹치지 않으므로 구조적으로 $0$ 이어야 하고, 실제로 $0$
이다.} 배선 확인이 겸사겸사 끝났다.

\section{씨앗 성분은 못 믿는다}

씨앗 상관 행렬도 냈다(판 $\leftrightarrow$ 날짜 통제 $+0.88$, KR
$\leftrightarrow$ CN $+0.86$, 무리 안 $\leftrightarrow$ 앱 $-0.60$, 유효
$3.02$). \textbf{쓰지 않는다.}

\begin{quote}
\textbf{씨앗 $8$ 점으로 $6\times6$ 상관을 재면 사실상 미결정이다.} 자유도가
$15$ 개인데 관측이 $8$ 개다. ``KR 과 CN 이 $+0.86$ 이다(둘 다 만화라
그럴듯하다)''는 \textbf{말이 되는 이야기가 붙는다는 것이 곧 잰 것은
아니다}(노트 $133$ 의 다른 얼굴). 표본 쪽은 $3{,}200$ 점이라 믿는다.
\end{quote}

\section{그래서 관문은 얼마나 센가}

\begin{center}
\begin{tabular}{lr}
\hline
 & 유효 자 수 \\
\hline
이름상 & $6$ \\
\textbf{표본 성분} & $\mathbf{4.73}$ \\
판 계열만 & $1.97$ \\
집 밖만 & $2.99$ \\
\hline
\end{tabular}
\end{center}

\textbf{집 밖 셋은 서로 완전히 독립이다}($2.99/3$, 상관 $|0.06|$ 아래).
판 계열 셋은 $\mathbf{1.97/3}$ --- 판과 날짜 통제 판이 사실상 하나다.

\begin{quote}
\textbf{그러나 자 수가 곧 힘은 아니다.} 노트 $613$ 이 봤듯 판 계열 셋은
이번 채택에 대해 \emph{셋 다 침묵했다}($t{=}0.7 \cdot 1.3 \cdot 0.2$).
\textbf{독립이어도 감도가 없으면 못 센다.} 실제로 이 채택을 결정한 것은
$4.73$ 중 \textbf{둘}이다.
\end{quote}

\section*{조항}

\begin{quote}
\textbf{관문의 힘은 자 수가 아니라 (독립 $\times$ 감도)다.} 이 판은 자를
여섯 두고 ``다 통과''를 강한 증거로 읽어 왔다. 독립은 $4.73$ 로 나쁘지
않은데, \textbf{그중 셋은 이 크기의 변경을 아예 못 잰다.} 자를 늘릴 때는
독립성과 감도를 \emph{같이} 봐야 한다 --- CN 만화는 독립성은 완벽했고
감도가 없어서 탈락했다(노트 $613$).
\end{quote}

\section*{남은 것}

\begin{itemize}
\item \textbf{감도를 자마다 미리 적어 두지 않았다.} $2$ SD 문턱은
노트 $613$ 이 냈지만, ``이 자로 잴 수 있는 최소 효과''를 후보 낼 때
자동으로 찍지는 않는다. \texttt{sideaudit} 에 안 넣었다.
\item \textbf{상관이 팔 짝마다 다를 수 있다.} 이 표는 \emph{채택 팔}
하나에서 잰 것이다. 가림 팔에서도 같은지 안 봤다.
\item \textbf{무리 안 판을 살릴 수 있나} --- 독립성은 좋은데($+0.11$)
감도가 없다($t{=}0.2$). 무리 문턱 $8$ 을 올리는 것이 후보인데
\textbf{잡음과 표본이 반대로 움직여} 어느 쪽이 이기는지 안 쟀다.
\end{itemize}

\end{document}
